\chapter*{Preface}
\addcontentsline{toc}{chapter}{Preface}
\markboth{Preface}{Preface}

\section*{What this book is about}

A robot lawnmower drives around a square lawn. We tell it how far to drive
and how much to turn, but it never does exactly that: wheels slip on wet
grass, the ground is uneven, the motors are imprecise. Its sensors are
noisy too. So where is it? This book answers that question in several ways,
each one a \emph{recursive Bayesian filter}: the grid filter, the Kalman
filter, the extended and the unscented Kalman filters, and the particle
filter. All of them are derived from one recursion, the Bayes filter, and
differ only in how they approximate it.

Every filter is also implemented and tested on the same simulated mower,
with the same data. This makes the comparisons honest. When a filter is
faster, more accurate or more honest about its own uncertainty than
another, the difference comes from the algorithm and not from the
experiment.

\section*{How to read it}

The book and the notebooks of the repository go together. The book gives
the theory: definitions, statements and complete proofs, with an
\emph{In words} box after each important result to say what it means. Each
chapter ends with a \emph{Practice} box that points to its notebook, where
the same ideas are implemented, run and plotted, and with exercises whose
solutions are given in full.

\begin{center}
\small
\begin{tabular}{@{}lll@{}}
\toprule
Chapter & Notebook (in \code{notebooks/}) & Code\\
\midrule
1 The state-space model & \code{00\_the\_mower\_world} & \code{world.py}\\
2 The Bayes filter & \code{01\_bayes\_filter} & \code{grid.py}\\
3 The Kalman filter & \code{02\_kalman\_filter} & \code{kf.py},
  \code{metrics.py}\\
4 The extended Kalman filter & \code{03\_extended\_kalman\_filter} &
  \code{ekf.py}\\
5 The unscented Kalman filter & \code{04\_unscented\_kalman\_filter} &
  \code{ukf.py}\\
6 The particle filter & \code{05\_particle\_filter} & \code{pf.py}\\
7 Choosing a filter & \code{06\_showdown} & all of them\\
\bottomrule
\end{tabular}
\end{center}

The boxes follow a fixed colour code: \textcolor{defblue}{blue} for
definitions, \textcolor{thmred}{red} for theorems,
\textcolor{propamber}{amber} for propositions and lemmas,
\textcolor{wordsgreen}{green} for the \emph{In words} intuition, grey for
practice, and a lighter green for the \emph{On the lawn} boxes, which say
what a result means for a real mower.

\section*{Prerequisites}

Linear algebra (symmetric positive definite matrices, the Schur
complement), multivariable calculus (Jacobians, Taylor expansions), and
probability at the level of a first course: densities, conditioning,
Bayes' rule, expectations and covariances, the multivariate Gaussian.
Chapter~\ref{ch:kalman} re-derives the Gaussian identities it needs. The
code is NumPy; no GPU is needed.

\section*{Building the book and the figures}

From the repository root, \code{make book} compiles \code{book/main.pdf}
and \code{make figures} regenerates every figure from the
\code{bayesian\_filters} package. The figures use the same random seeds as
the notebooks, so the numbers quoted in the text are the numbers printed
by the notebooks.

\section*{Where to go next}

Filtering estimates the \emph{current} state. Modern localisation and
mapping (SLAM) systems instead estimate a whole trajectory at once, by
nonlinear least squares over a factor graph. Chapter~\ref{ch:choosing}
ends with a short bridge to that point of view, and the companion course
\url{https://github.com/CL-BZH/vslam-course} develops it in full for
visual SLAM on the same robot mower.
